% Fentes d'Young : différence de marche et franges d'interférences (L'essentiel, chapitre 22)
\begin{tikzpicture}[x=1cm,y=1cm]
\draw[otRed!80!black,line width=1.1pt] (-1.6,0) -- (0,0);
\node[pcLegende,anchor=south] at (-1.0,0.05) {laser};
% plaque et fentes
\fill[otInk!70] (0,0.42) rectangle (0.12,1.4) (0,-0.32) rectangle (0.12,0.32) (0,-1.4) rectangle (0.12,-0.42);
\node[pcLegende,anchor=west] at (0.18,0.55) {$S_1$};
\node[pcLegende,anchor=west] at (0.18,-0.55) {$S_2$};
\draw[otMuted,line width=0.4pt,{Stealth[length=1.1mm]}-{Stealth[length=1.1mm]}] (-0.2,-0.37) -- (-0.2,0.37);
\node[pcLegende,anchor=east] at (-0.25,0) {$b$};
% rayons vers M
\fill[otInk] (6.0,1.05) circle (0.05);
\draw[otNavy,line width=0.6pt] (0.12,0.37) -- (6.0,1.05);
\draw[otNavy,line width=0.6pt] (0.12,-0.37) -- (6.0,1.05);
\node[pcLegende,anchor=south east] at (5.95,1.08) {$M$};
\draw[otMuted,line width=0.4pt,dash pattern=on 2pt off 1.5pt] (0.12,0) -- (6.0,0);
% écran et franges
\fill[otMuted!25] (6.0,-1.6) rectangle (6.1,1.6);
\foreach \k in {-3,...,3}{\fill[otRed!80] (6.1,{0.5*\k-0.11}) rectangle (6.35,{0.5*\k+0.11});}
\foreach \k in {-3,...,2}{\fill[otRed!20] (6.1,{0.5*\k+0.11}) rectangle (6.35,{0.5*\k+0.39});}
\draw[otMuted,line width=0.4pt,{Stealth[length=1.1mm]}-{Stealth[length=1.1mm]}] (6.55,0) -- (6.55,0.5) node[pos=0.5,pcLegende,anchor=west] {$i$};
\draw[otMuted,line width=0.4pt,-{Stealth[length=1.2mm]}] (6.45,-1.65) -- (6.45,1.75) node[pcLegende,anchor=south] {$x$};
\draw[otMuted,line width=0.4pt,{Stealth[length=1.2mm]}-{Stealth[length=1.2mm]}] (0.12,-1.75) -- (6.0,-1.75) node[pos=0.5,pcLegende,anchor=north] {$D$};
\node[pcLegende,anchor=north west,align=left] at (1.3,-0.25) {différence de marche en $M$ :\\ $\delta=S_2M-S_1M\approx\dfrac{b\,x}{D}$};
\end{tikzpicture}
\par\smallskip\legende{Les deux fentes, éclairées par le même laser, se comportent comme deux sources en phase ; sur l'écran alternent franges brillantes et sombres.}
